\section{Evaluation}
\label{sec:evaluation}

\subsection{Concept Questions}

\draftnote{Describe the declarative completion format, train/test question splits, concept and neighboring-domain sets, and option scoring. Define the primary continuous gold-answer likelihood statistic and treat multiple-choice accuracy as descriptive. Explain why generation and letter parsing are inappropriate for these base models.}

\subsection{Held-Out Text Likelihood}

\draftnote{Define per-token negative log-likelihood on identical concept chunks and length-matched control chunks. Give the difference-in-differences estimand explicitly: the concept-chunk model difference minus the control-chunk model difference. Report sample sizes and pairing by chunk ID.}

\subsection{Parameter-Space Comparison}

\draftnote{Define relative edit size, cosine similarity between the erasure update and never-training displacement, progress along the never-trained direction, and residual distance. State which tensors are included and separately report the embedding-only comparison where useful.}

\subsection{Statistical Analysis}

\draftnote{Specify paired tests, permutation procedures, effect sizes, and confidence intervals. Distinguish item-level uncertainty from unmeasured training-seed uncertainty. Document which data selected EMBER's delta and reserve held-out splits for final claims.}

